
\documentclass[aspectratio=169,11pt]{beamer}

% ---------- Packages ----------
\usepackage[utf8]{inputenc}
\usepackage[T1]{fontenc}
\usepackage{lmodern}
\usepackage{amsmath,amssymb,bm}
\usepackage{booktabs}
\usepackage{array}
\usepackage{hyperref}
\usepackage{xcolor}
\usepackage{multicol}
\usepackage{microtype}

% ---------- Theme ----------
\usetheme{Madrid}
\usecolortheme{default}
\definecolor{WLUBlue}{HTML}{0047AB}
\definecolor{WLUGray}{HTML}{444444}
\definecolor{WLULight}{HTML}{EEF4FF}
\definecolor{WLUGold}{HTML}{B9985A}

\setbeamercolor{structure}{fg=WLUBlue}
\setbeamercolor{title}{fg=white,bg=WLUBlue}
\setbeamercolor{frametitle}{fg=white,bg=WLUBlue}
\setbeamercolor{block title}{fg=white,bg=WLUBlue}
\setbeamercolor{block body}{bg=WLULight,fg=black}
\setbeamercolor{alertblock title}{fg=white,bg=WLUGold}
\setbeamercolor{alertblock body}{bg=WLUGold!12,fg=black}
\setbeamercolor{exampleblock title}{fg=white,bg=WLUGray}
\setbeamercolor{exampleblock body}{bg=black!4,fg=black}

\setbeamertemplate{navigation symbols}{}
\setbeamertemplate{footline}{%
  \hfill\scriptsize Module 1: Economic Questions, Prediction, and Causality \quad Slide \insertframenumber\hspace{1em}\vspace{0.6em}
}

\hypersetup{
  colorlinks=true,
  linkcolor=WLUBlue,
  urlcolor=WLUBlue,
  citecolor=WLUBlue
}

% ---------- Helpers ----------
\newcommand{\core}[1]{\begin{block}{Core intuition}#1\end{block}}
\newcommand{\grad}[1]{\begin{alertblock}{Technical layer}#1\end{alertblock}}
\newcommand{\econ}[1]{\begin{exampleblock}{Economic interpretation}#1\end{exampleblock}}
\newcommand{\smallcite}[1]{{\scriptsize \textcolor{WLUGray}{#1}}}

\title[ML in Economics]{Module 1: Economic Questions, Prediction, and Causality}
\subtitle{Machine Learning in Economics}
\author{EC 410K / EC 610I}
\institute{Wilfrid Laurier University}
\date{Spring 2026}

\begin{document}

% =========================================================
\begin{frame}
  \titlepage
  \vspace{-0.5em}
  \begin{center}
  \small
  This module builds the conceptual bridge between econometrics, machine learning, and policy-relevant empirical work.
  \end{center}
\end{frame}

% =========================================================
\begin{frame}{Why this module matters}
\core{
Before choosing Lasso, random forests, neural networks, instrumental variables, regression discontinuity, or difference-in-differences, we must decide what question we are answering.
}
\begin{itemize}
  \item Machine learning is excellent at discovering predictive patterns.
  \item Economics often asks causal and policy questions.
  \item The empirical strategy should follow the economic question, not the other way around.
\end{itemize}
\econ{
A model that predicts poverty well may still fail to identify whether a transfer program reduces poverty.
}
\end{frame}

% =========================================================
\begin{frame}{Learning goals}
By the end of Module 1, students should be able to:
\begin{enumerate}
  \item Distinguish prediction questions from causal questions.
  \item Translate an economic question into an estimand.
  \item Explain why endogeneity and selection bias create causal problems.
  \item Identify the main logic behind common identification strategies.
  \item Explain how machine learning can support, but not replace, identification.
  \item Read replication papers with a clear map: question, data, design, estimand, threats, and possible ML extension.
\end{enumerate}
\end{frame}

% =========================================================
\begin{frame}{How the module connects to the course}
\begin{columns}
\begin{column}{0.48\textwidth}
\begin{block}{Module 1 foundation}
\begin{itemize}
  \item Economic questions
  \item Prediction versus causality
  \item Endogeneity
  \item Selection
  \item Identification
\end{itemize}
\end{block}
\end{column}
\begin{column}{0.48\textwidth}
\begin{alertblock}{Later modules build on this}
\begin{itemize}
  \item Supervised learning
  \item Regularization
  \item Causal ML
  \item Trees and interpretability
  \item Forecasting
  \item NLP and generative AI
  \item Replication and project work
\end{itemize}
\end{alertblock}
\end{column}
\end{columns}
\end{frame}

% =========================================================
\begin{frame}{The central distinction}
\begin{center}
\Large
Prediction asks: \quad ``What will happen?''

\vspace{1em}

Causal inference asks: \quad ``What would happen if we changed something?''
\end{center}

\vspace{1em}
\begin{block}{A simple rule}
If the question contains a policy, treatment, intervention, rule, threshold, exposure, institution, or historical shock, it is probably a causal question.
\end{block}
\end{frame}

% =========================================================
\begin{frame}{Prediction questions}
\core{
Prediction aims to estimate an unknown outcome for an observation.
}
\[
\widehat{Y}_i = \widehat{f}(X_i)
\]
\begin{itemize}
  \item Who is likely to default on a loan?
  \item Which counties will experience employment decline?
  \item Which Medicare drug plans will attract subsidized enrollees?
  \item Which regions are likely to have low income today?
\end{itemize}
\grad{
The target is a prediction function $\widehat f$, evaluated using out-of-sample performance such as RMSE, MAE, accuracy, AUC, or log loss.
}
\end{frame}

% =========================================================
\begin{frame}{Causal questions}
\core{
Causal inference asks about a counterfactual: what would have happened to the same unit under a different treatment or policy?
}
\[
\tau_i = Y_i(1)-Y_i(0)
\]
\begin{itemize}
  \item Did colonial institutions affect modern income?
  \item Did exposure to slave trades affect current economic performance?
  \item Did stronger drunk-driving sanctions reduce recidivism?
  \item Did import competition from China reduce U.S. manufacturing employment?
\end{itemize}
\grad{
The target is usually an estimand such as ATE, ATT, LATE, CATE, or a policy value function.
}
\end{frame}

% =========================================================
\begin{frame}{Prediction versus causality: same data, different question}
Suppose we observe workers with earnings $Y$, training status $D$, and characteristics $X$.
\begin{columns}
\begin{column}{0.48\textwidth}
\begin{block}{Prediction}
\[
\widehat{Y} = \widehat f(D,X)
\]
Goal: predict earnings accurately.
\begin{itemize}
  \item Use all predictive variables.
  \item Care about test error.
  \item Coefficients need not be causal.
\end{itemize}
\end{block}
\end{column}
\begin{column}{0.48\textwidth}
\begin{alertblock}{Causality}
\[
\tau = E[Y(1)-Y(0)]
\]
Goal: estimate effect of training.
\begin{itemize}
  \item Need a credible comparison group.
  \item Need assumptions or design.
  \item Predictive variables may be confounders, mediators, or colliders.
\end{itemize}
\end{alertblock}
\end{column}
\end{columns}
\end{frame}

% =========================================================
\begin{frame}{A prediction model can be useful and non-causal}
\begin{block}{Example}
A random forest predicts unemployment using age, education, industry, location, and prior earnings.
\end{block}
\begin{itemize}
  \item It may identify workers at risk.
  \item It may help target outreach.
  \item But it does not tell us whether a training program works.
\end{itemize}
\econ{
Prediction can improve policy administration even when it does not estimate causal effects. For example, a government agency may use prediction to screen who is at risk, then evaluate the treatment using a randomized or quasi-experimental design.
}
\end{frame}

% =========================================================
\begin{frame}{Why causal inference is hard}
For each unit $i$, we would like to observe both:
\[
Y_i(1) \quad \text{and} \quad Y_i(0)
\]
But we only observe one:
\[
Y_i = D_iY_i(1) + (1-D_i)Y_i(0)
\]
\begin{block}{Fundamental problem of causal inference}
We never observe the same unit at the same time both treated and untreated.
\end{block}
\end{frame}

% =========================================================
\begin{frame}{The econometric model is not the identification strategy}
A regression equation may look causal:
\[
Y_i = \alpha + \theta D_i + X_i'\beta + \varepsilon_i
\]
But the coefficient $\theta$ is causal only if the variation in $D_i$ is as-good-as-random after conditioning on $X_i$.
\begin{alertblock}{Key warning}
Adding more controls, using nonlinear ML, or achieving high $R^2$ does not automatically solve endogeneity.
\end{alertblock}
\end{frame}

% =========================================================
\begin{frame}{Selection bias: the basic problem}
Observed difference in outcomes:
\[
E[Y|D=1]-E[Y|D=0]
\]
can be decomposed as:
\[
\underbrace{E[Y(1)-Y(0)|D=1]}_{\text{treatment effect on treated}}
+
\underbrace{E[Y(0)|D=1]-E[Y(0)|D=0]}_{\text{selection bias}}
\]
\core{
Selection bias arises when treated and untreated units would have differed even without treatment.
}
\end{frame}

% =========================================================
\begin{frame}{Selection bias in economic examples}
\begin{itemize}
  \item Students who receive tutoring may be more motivated.
  \item Counties exposed to trade shocks may have different industrial structures before the shock.
  \item Countries with stronger institutions may differ in geography, disease environment, or colonial history.
  \item Medicare plans near subsidy thresholds may differ strategically from plans far away.
\end{itemize}
\econ{
The empirical task is not just to compare groups; it is to construct a comparison that approximates the missing counterfactual.
}
\end{frame}

% =========================================================
\begin{frame}{Endogeneity: three common channels}
\begin{enumerate}
  \item \textbf{Omitted variables:} $D$ is correlated with unobserved determinants of $Y$.
  \item \textbf{Reverse causality:} $Y$ affects $D$, not only the other way around.
  \item \textbf{Measurement error:} the observed treatment or key regressor is noisy.
\end{enumerate}
\begin{alertblock}{Economic implication}
A flexible algorithm can estimate the wrong relationship very accurately if the source of variation is not credible.
\end{alertblock}
\end{frame}

% =========================================================
\begin{frame}{Identification: the credibility problem}
\core{
Identification is the argument that the data contain variation capable of answering the causal question.
}
\begin{itemize}
  \item What is the source of variation in treatment?
  \item Why is it plausibly unrelated to potential outcomes?
  \item What is the relevant counterfactual group?
  \item What assumptions are required?
  \item What evidence can support or challenge those assumptions?
\end{itemize}
\grad{
Estimation is about computing a number. Identification is about whether that number has the meaning we assign to it.
}
\end{frame}

% =========================================================
\begin{frame}{Common causal designs}
\begin{center}
\begin{tabular}{p{0.25\textwidth}p{0.32\textwidth}p{0.30\textwidth}}
\toprule
Design & Source of variation & Key assumption \\
\midrule
Randomized experiment & Random assignment & Randomization worked \\
Instrumental variables & Instrument shifts treatment & Exclusion and relevance \\
Regression discontinuity & Threshold rule & Continuity around cutoff \\
Difference-in-differences & Before-after and treated-control & Parallel trends \\
Propensity score methods & Observed covariates & Unconfoundedness and overlap \\
Shift-share design & Differential exposure to aggregate shocks & Valid shares and shocks \\
\bottomrule
\end{tabular}
\end{center}
\end{frame}

% =========================================================
\begin{frame}{Where machine learning enters}
Machine learning can help with:
\begin{itemize}
  \item prediction of outcomes,
  \item flexible control functions,
  \item high-dimensional covariate adjustment,
  \item variable selection,
  \item propensity score estimation,
  \item heterogeneous treatment effects,
  \item policy targeting,
  \item text-as-data measurement,
  \item robustness and diagnostics.
\end{itemize}
\begin{alertblock}{But}
Machine learning does not itself create exogenous variation.
\end{alertblock}
\end{frame}

% =========================================================
\begin{frame}{Two roles for ML in economics}
\begin{columns}
\begin{column}{0.48\textwidth}
\begin{block}{ML as the main object}
\begin{itemize}
  \item Forecast inflation.
  \item Predict loan default.
  \item Classify recession risk.
  \item Extract sentiment from text.
\end{itemize}
Evaluation: prediction quality and usefulness.
\end{block}
\end{column}
\begin{column}{0.48\textwidth}
\begin{alertblock}{ML as a supporting tool}
\begin{itemize}
  \item Estimate nuisance functions.
  \item Select controls.
  \item Explore heterogeneity.
  \item Improve policy assignment.
\end{itemize}
Evaluation: causal validity plus statistical performance.
\end{alertblock}
\end{column}
\end{columns}
\end{frame}

% =========================================================
\begin{frame}{Terminology map}
\begin{center}
\begin{tabular}{p{0.25\textwidth}p{0.33\textwidth}p{0.32\textwidth}}
\toprule
Term & Prediction meaning & Causal meaning \\
\midrule
Target & $Y$ to forecast & Effect or counterfactual \\
Model & $\widehat f(X)$ & Estimand plus design \\
Error & Test loss & Bias, variance, invalid inference \\
Controls & Predictive features & Confounders to adjust for \\
Validation & Holdout or CV & Balance, pre-trends, placebo tests, sensitivity \\
Interpretability & Feature importance & Economic mechanism and credible assumptions \\
\bottomrule
\end{tabular}
\end{center}
\end{frame}

% =========================================================
\begin{frame}{Estimands: what exactly are we trying to learn?}
\begin{block}{Average treatment effect}
\[
ATE = E[Y(1)-Y(0)]
\]
\end{block}
\begin{block}{Average treatment effect on treated}
\[
ATT = E[Y(1)-Y(0)|D=1]
\]
\end{block}
\begin{block}{Conditional average treatment effect}
\[
CATE(x)=E[Y(1)-Y(0)|X=x]
\]
\end{block}
\econ{
Policy often cares about CATE: not only whether a program works on average, but for whom it works best.
}
\end{frame}

% =========================================================
\begin{frame}{From research question to estimand}
\begin{center}
\begin{tabular}{p{0.30\textwidth}p{0.30\textwidth}p{0.30\textwidth}}
\toprule
Question & Candidate estimand & Example design \\
\midrule
Did a policy work? & ATE or ATT & RCT, DiD, RD \\
Who benefits most? & CATE & Causal forest, subgroup analysis \\
What happens near a rule cutoff? & Local treatment effect & RD \\
What is the effect for compliers? & LATE & IV \\
How should we target a policy? & Policy value & Policy learning \\
\bottomrule
\end{tabular}
\end{center}
\end{frame}

% =========================================================
\begin{frame}{A minimal empirical workflow}
\begin{enumerate}
  \item State the economic question.
  \item Define units, outcome, treatment, and timing.
  \item Decide whether the question is predictive, causal, or both.
  \item Specify the estimand.
  \item Identify the source of variation.
  \item Choose the estimator.
  \item Validate the design and model.
  \item Interpret results economically.
  \item Discuss limitations and external validity.
\end{enumerate}
\end{frame}

% =========================================================
\begin{frame}{What makes an economic question good?}
A strong empirical question is:
\begin{itemize}
  \item \textbf{specific:} clear units, period, outcome, and treatment;
  \item \textbf{economically meaningful:} linked to theory, policy, or behavior;
  \item \textbf{empirically feasible:} data can measure the key objects;
  \item \textbf{identifiable:} credible comparison exists;
  \item \textbf{interpretable:} results can be explained to economists and policymakers.
\end{itemize}
\begin{block}{Bad question}
Does technology affect workers?
\end{block}
\begin{alertblock}{Better question}
How did rising Chinese import competition affect manufacturing employment in U.S. commuting zones from 1990 to 2007?
\end{alertblock}
\end{frame}

% =========================================================
\begin{frame}{Reading empirical papers: the six-part checklist}
For every paper in the course replication list, ask:
\begin{enumerate}
  \item What is the economic question?
  \item What is the outcome?
  \item What is the treatment, exposure, or policy variable?
  \item What is the comparison group or source of variation?
  \item What could make the estimate biased?
  \item How could machine learning extend the analysis?
\end{enumerate}
\end{frame}

% =========================================================
\begin{frame}{Replication papers used throughout this module}
This module uses examples from papers students may later replicate and extend:
\begin{enumerate}
  \item Acemoglu, Johnson, and Robinson (2001): colonial origins and institutions.
  \item Nunn (2008): slave trades and current economic performance.
  \item Goldin, Lurie, and McCubbin (2016): Medicare Part D subsidy design.
  \item Hansen (2015): punishment and deterrence in drunk driving.
  \item Anderson (2017): college athletic success and institutional outcomes.
  \item Autor, Dorn, and Hanson (2013): China import competition and local labor markets.
\end{enumerate}
\smallcite{These are listed in the course outline as replication and ML-extension papers.}
\end{frame}

% =========================================================
\begin{frame}{Paper example 1: AJR (2001)}
\begin{block}{Question}
Do colonial-era institutions help explain current differences in economic development?
\end{block}
\begin{itemize}
  \item Outcome: modern income per capita.
  \item Key explanatory variable: institutional quality.
  \item Empirical challenge: institutions are endogenous.
  \item Strategy: use historical settler mortality as an instrument for institutions.
\end{itemize}
\smallcite{Acemoglu, Johnson, and Robinson (2001), American Economic Review.}
\end{frame}

% =========================================================
\begin{frame}{AJR as a prediction question}
A pure prediction version might ask:
\begin{center}
\Large
Can we predict current GDP per capita using geography, colonial history, institutions, disease environment, and regional variables?
\end{center}
\vspace{0.5em}
Possible ML tools:
\begin{itemize}
  \item random forests or boosting to predict income,
  \item Lasso to select strong predictors,
  \item cross-validation to evaluate predictive performance.
\end{itemize}
\begin{alertblock}{Limitation}
High predictive performance would not prove that institutions caused income differences.
\end{alertblock}
\end{frame}

% =========================================================
\begin{frame}{AJR as a causal question}
The causal version asks:
\begin{center}
\Large
What is the effect of institutions on long-run economic development?
\end{center}
\[
\log GDP_i = \alpha + \theta Institutions_i + X_i'\beta + u_i
\]
\begin{itemize}
  \item $Institutions_i$ may be correlated with geography, colonial strategy, and unobserved historical factors.
  \item The IV strategy seeks variation in institutions caused by settler mortality.
  \item The central debate is whether the instrument affects income only through institutions.
\end{itemize}
\end{frame}

% =========================================================
\begin{frame}{AJR: what can ML add?}
\begin{itemize}
  \item Predict first-stage institutional quality using high-dimensional historical and geographic variables.
  \item Use ML to test robustness to many alternative controls.
  \item Explore heterogeneity: are institutional effects different by continent, resource endowment, or colonial power?
  \item Use explainability tools to compare predictive importance with causal interpretation.
\end{itemize}
\begin{alertblock}{Critical caution}
ML can strengthen prediction and robustness checks, but the IV exclusion restriction remains an economic and historical argument.
\end{alertblock}
\end{frame}

% =========================================================
\begin{frame}{Paper example 2: Nunn (2008)}
\begin{block}{Question}
Did historical slave trades have persistent effects on current economic performance in Africa?
\end{block}
\begin{itemize}
  \item Outcome: modern income or economic performance.
  \item Treatment/exposure: intensity of slave exports.
  \item Empirical challenge: regions with greater slave exports may differ systematically.
  \item Identification relies on historical and geographic variation linked to slave trade intensity.
\end{itemize}
\smallcite{Nunn (2008), Quarterly Journal of Economics.}
\end{frame}

% =========================================================
\begin{frame}{Nunn: prediction versus causality}
\begin{columns}
\begin{column}{0.48\textwidth}
\begin{block}{Prediction question}
Can historical, geographic, and institutional variables predict current income across countries or ethnic groups?
\end{block}
\end{column}
\begin{column}{0.48\textwidth}
\begin{alertblock}{Causal question}
What was the effect of slave trade exposure on long-run development outcomes?
\end{alertblock}
\end{column}
\end{columns}
\vspace{1em}
\econ{
The causal claim requires an argument that the measured exposure is not merely proxying for pre-existing development, geography, conflict, or other long-run factors.
}
\end{frame}

% =========================================================
\begin{frame}{Nunn: ML extensions}
Possible student project extensions:
\begin{itemize}
  \item Use regularization to assess which historical controls matter most.
  \item Use tree-based models to predict current economic performance and compare with linear specifications.
  \item Use causal forests to explore whether persistence differs by colonial history, geography, or precolonial institutions.
  \item Use sensitivity analysis to assess robustness to alternative control sets.
\end{itemize}
\grad{
A useful project should separate ``ML improves prediction'' from ``ML improves the credibility or precision of a causal estimate.''
}
\end{frame}

% =========================================================
\begin{frame}{Historical persistence papers: common identification issues}
AJR and Nunn share broad challenges:
\begin{itemize}
  \item Many causes of current income are historically intertwined.
  \item Treatment variables are measured with historical uncertainty.
  \item Geographic, institutional, and cultural mechanisms may overlap.
  \item External validity and interpretation require careful historical context.
\end{itemize}
\begin{alertblock}{Lesson for ML}
High-dimensional data can help organize evidence, but historical causal claims still require explicit identification assumptions.
\end{alertblock}
\end{frame}

% =========================================================
\begin{frame}{Paper example 3: Goldin, Lurie, and McCubbin (2016)}
\begin{block}{Question}
Do insurers strategically respond to the design of the Medicare Part D Low-Income Subsidy?
\end{block}
\begin{itemize}
  \item Setting: Medicare Part D prescription drug plans.
  \item Policy feature: subsidy rules create incentives around benchmark premiums.
  \item Outcome: plan pricing, enrollment, or behavior around subsidy eligibility.
  \item Empirical challenge: distinguishing strategic response from ordinary plan differences.
\end{itemize}
\smallcite{Goldin, Lurie, and McCubbin (2016), American Economic Review.}
\end{frame}

% =========================================================
\begin{frame}{Goldin et al.: prediction versus causality}
\begin{block}{Prediction version}
Which plans are likely to attract low-income subsidy enrollees?
\end{block}
\begin{alertblock}{Causal/policy version}
Did the subsidy design induce insurers to manipulate or strategically set plan premiums?
\end{alertblock}
\begin{itemize}
  \item Prediction can help identify high-risk plans or markets.
  \item Causality requires understanding incentives created by the policy formula.
  \item Institutional knowledge is central: the rule itself shapes behavior.
\end{itemize}
\end{frame}

% =========================================================
\begin{frame}{Goldin et al.: ML extensions}
Possible ML extensions:
\begin{itemize}
  \item Classification model: predict which plans will locate near benchmark thresholds.
  \item Tree-based heterogeneity: examine whether strategic behavior differs by market concentration, plan type, or region.
  \item Anomaly detection: flag unusual premium-setting patterns.
  \item Counterfactual simulation: compare actual pricing to predicted pricing absent the incentive.
\end{itemize}
\begin{alertblock}{Identification caution}
A model that detects bunching or unusual pricing is suggestive; causal interpretation depends on the policy rule and counterfactual behavior.
\end{alertblock}
\end{frame}

% =========================================================
\begin{frame}{Paper example 4: Hansen (2015)}
\begin{block}{Question}
Does punishment deter drunk driving?
\end{block}
\begin{itemize}
  \item Setting: drunk-driving cases.
  \item Key institutional feature: penalties increase sharply at blood alcohol content thresholds.
  \item Outcome: future drunk-driving behavior, recidivism, or related outcomes.
  \item Design logic: compare individuals just below and just above the threshold.
\end{itemize}
\smallcite{Hansen (2015), American Economic Review.}
\end{frame}

% =========================================================
\begin{frame}{Hansen: regression discontinuity intuition}
Regression discontinuity uses a cutoff rule:
\[
D_i = 1\{RunningVariable_i \geq c\}
\]
\begin{itemize}
  \item Units just below and just above the cutoff are assumed similar.
  \item A discontinuity in outcomes at the cutoff is attributed to treatment.
  \item The estimand is local: it applies near the cutoff.
\end{itemize}
\core{
RD is compelling because the comparison group is constructed by the policy rule itself.
}
\end{frame}

% =========================================================
\begin{frame}{Hansen: prediction versus causality}
\begin{block}{Prediction}
Can we predict which drivers will reoffend?
\end{block}
\begin{alertblock}{Causality}
Does crossing the legal punishment threshold reduce future drunk-driving behavior?
\end{alertblock}
\econ{
Prediction might use age, prior record, location, and BAC. Causal inference focuses on whether the threshold creates quasi-random treatment assignment near the cutoff.
}
\end{frame}

% =========================================================
\begin{frame}{Hansen: ML extensions}
Possible ML extensions:
\begin{itemize}
  \item Use ML to predict recidivism risk and compare with RD estimates.
  \item Estimate heterogeneous treatment effects near the cutoff.
  \item Use flexible functional forms for the running variable while preserving RD logic.
  \item Study whether deterrence differs by age, prior record, or local enforcement intensity.
\end{itemize}
\grad{
For RD, flexible ML must be used carefully. Overfitting near the cutoff can distort the estimated discontinuity.
}
\end{frame}

% =========================================================
\begin{frame}{Paper example 5: Anderson (2017)}
\begin{block}{Question}
Do successful college athletic programs generate broader institutional benefits?
\end{block}
\begin{itemize}
  \item Outcomes: applications, enrollment, donations, academic indicators, or institutional outcomes.
  \item Treatment: athletic success.
  \item Challenge: successful sports programs may differ from other institutions in many ways.
  \item Design: propensity score approach to improve comparability.
\end{itemize}
\smallcite{Anderson (2017), Review of Economics and Statistics.}
\end{frame}

% =========================================================
\begin{frame}{Propensity score idea}
The propensity score is:
\[
e(X_i)=P(D_i=1|X_i)
\]
\begin{block}{Purpose}
Compare treated and untreated units with similar probability of receiving treatment.
\end{block}
\begin{itemize}
  \item Balances observed covariates.
  \item Reduces dimensionality of matching.
  \item Requires no unobserved confounding after conditioning on $X$.
\end{itemize}
\end{frame}

% =========================================================
\begin{frame}{Anderson: prediction versus causality}
\begin{block}{Prediction}
Can we predict which colleges will receive more applications or donations?
\end{block}
\begin{alertblock}{Causality}
Did athletic success cause increases in applications, enrollment, or donations?
\end{alertblock}
\econ{
Predictive success may reflect brand, selectivity, size, conference membership, alumni networks, or location. The causal design tries to compare similar institutions differing in athletic success.
}
\end{frame}

% =========================================================
\begin{frame}{Anderson: ML extensions}
Possible ML extensions:
\begin{itemize}
  \item Estimate propensity scores using flexible methods such as random forests, boosting, or Lasso.
  \item Check covariate balance before and after matching or weighting.
  \item Estimate heterogeneous effects by institution type, conference, baseline selectivity, or media exposure.
  \item Compare conventional propensity score estimates with doubly robust ML estimates.
\end{itemize}
\grad{
ML propensity scores can reduce misspecification, but they do not solve unobserved confounding.
}
\end{frame}

% =========================================================
\begin{frame}{Paper example 6: Autor, Dorn, and Hanson (2013)}
\begin{block}{Question}
How did rising import competition from China affect U.S. local labor markets?
\end{block}
\begin{itemize}
  \item Unit: local labor market, often commuting zone.
  \item Treatment/exposure: import competition exposure.
  \item Outcomes: manufacturing employment, wages, labor-force participation, transfers, and other local outcomes.
  \item Strategy: exploit variation in local industry composition and changes in Chinese import growth.
\end{itemize}
\smallcite{Autor, Dorn, and Hanson (2013), American Economic Review.}
\end{frame}

% =========================================================
\begin{frame}{Autor-Dorn-Hanson: shift-share intuition}
Local areas differ in pre-existing industry shares.
\[
Exposure_{cz} = \sum_j Share_{cz,j,0}\times ImportGrowth_j
\]
\begin{itemize}
  \item Areas initially specialized in industries with large import growth become more exposed.
  \item The strategy compares local labor markets with different predicted exposure.
  \item A major concern is whether initial shares are correlated with other local trends.
\end{itemize}
\end{frame}

% =========================================================
\begin{frame}{Autor-Dorn-Hanson: prediction versus causality}
\begin{block}{Prediction}
Can we predict which commuting zones will lose manufacturing jobs?
\end{block}
\begin{alertblock}{Causality}
Did Chinese import competition cause local labor-market decline?
\end{alertblock}
\begin{itemize}
  \item Prediction may use demographics, industry mix, education, unionization, regional trends, and trade exposure.
  \item Causal interpretation depends on the validity of the exposure measure and instrument.
\end{itemize}
\end{frame}

% =========================================================
\begin{frame}{Autor-Dorn-Hanson: ML extensions}
Possible ML extensions:
\begin{itemize}
  \item Predict local labor-market vulnerability using pre-shock covariates.
  \item Use regularization to select controls from rich local data.
  \item Estimate heterogeneous effects across regions, skill groups, or baseline industry structures.
  \item Compare policy targeting rules based on predicted vulnerability versus estimated causal harm.
\end{itemize}
\grad{
A useful extension asks whether ML improves prediction, identifies heterogeneity, or changes substantive policy conclusions.
}
\end{frame}

% =========================================================
\begin{frame}{Cross-paper comparison}
\begin{center}
\tiny
\begin{tabular}{p{0.17\textwidth}p{0.19\textwidth}p{0.22\textwidth}p{0.31\textwidth}}
\toprule
Paper & Main design idea & Main threat & ML opportunity \\
\midrule
AJR & IV using historical mortality & Exclusion restriction & Robust controls, heterogeneity \\
Nunn & Historical exposure variation & Omitted historical factors & High-dimensional controls, heterogeneity \\
Goldin et al. & Policy incentives and strategic response & Counterfactual pricing & Classification, anomaly detection \\
Hansen & RD at punishment threshold & Manipulation or local misspecification & Heterogeneity near cutoff \\
Anderson & Propensity score design & Unobserved confounding & ML propensity scores, DR methods \\
ADH & Shift-share exposure and IV logic & Correlated local trends & Vulnerability prediction, heterogeneous effects \\
\bottomrule
\end{tabular}
\end{center}
\end{frame}

% =========================================================
\begin{frame}{A recurring pattern across the papers}
\begin{enumerate}
  \item Start with an economically important question.
  \item Find a setting where treatment or exposure varies.
  \item Argue why the variation is informative about causality.
  \item Estimate the effect.
  \item Conduct robustness checks.
  \item Interpret the magnitude and mechanism.
  \item Discuss external validity.
\end{enumerate}
\begin{block}{Role of ML in your project}
ML should extend this logic, not replace it.
\end{block}
\end{frame}

% =========================================================
\begin{frame}{Prediction can be policy-relevant}
Prediction is not inferior to causality. It answers a different question.
\begin{itemize}
  \item Forecast recessions or inflation.
  \item Predict which households are likely to need support.
  \item Predict which regions are vulnerable to trade shocks.
  \item Predict which students are at risk of dropping out.
\end{itemize}
\econ{
A prediction model can improve resource allocation even if it does not estimate causal effects.
}
\end{frame}

% =========================================================
\begin{frame}{But prediction can create policy mistakes}
\begin{alertblock}{Mistake}
Use a predictive correlation as if it were a causal effect.
\end{alertblock}
\begin{itemize}
  \item If high school GPA predicts college success, forcing GPA to rise may not cause success.
  \item If regions with strong institutions are richer, simply copying institutional labels may not reproduce historical conditions.
  \item If high-risk drivers are predicted to reoffend, punishing them more severely may not reduce reoffense.
\end{itemize}
\core{
Prediction helps identify risk. Causal inference helps evaluate interventions.
}
\end{frame}

% =========================================================
\begin{frame}{Policy targeting needs both prediction and causality}
A policymaker may ask:
\begin{enumerate}
  \item Who is at risk?
  \item Which intervention works?
  \item For whom does it work best?
  \item What is the cost-benefit tradeoff?
\end{enumerate}
\[
\text{Target if } \widehat{\tau}(X_i)\times Benefit_i > Cost_i
\]
\grad{
Policy targeting requires estimating heterogeneous treatment effects, not just predicting untreated outcomes.
}
\end{frame}

% =========================================================
\begin{frame}{Good controls, bad controls, and ML features}
\begin{block}{Good control}
A pre-treatment variable that affects treatment and outcome.
\end{block}
\begin{alertblock}{Bad control}
A variable affected by the treatment, or a collider created by conditioning.
\end{alertblock}
\begin{itemize}
  \item ML encourages adding many features.
  \item Economics requires knowing when a feature belongs in the model.
  \item Timing matters: pre-treatment variables are generally safer than post-treatment variables.
\end{itemize}
\end{frame}

% =========================================================
\begin{frame}{Example: bad controls in a college athletics study}
Suppose athletic success affects applications, which affect selectivity, which affects donations.
\[
AthleticSuccess \rightarrow Applications \rightarrow Selectivity \rightarrow Donations
\]
If we control for applications or selectivity when estimating the effect on donations, we may block part of the causal pathway.
\begin{alertblock}{Lesson}
A variable can be highly predictive and still be a bad control for a causal question.
\end{alertblock}
\end{frame}

% =========================================================
\begin{frame}{Train-test splits are not enough for causal claims}
\begin{block}{Prediction validation}
Train on one sample and test on another.
\end{block}
\begin{alertblock}{Causal validation}
Ask whether the treatment variation is credible.
\end{alertblock}
\begin{itemize}
  \item Randomization checks.
  \item Covariate balance.
  \item Placebo outcomes.
  \item Pre-trend tests.
  \item Density or manipulation tests around cutoffs.
  \item Sensitivity to bandwidths, controls, and samples.
\end{itemize}
\end{frame}

% =========================================================
\begin{frame}{Overfitting versus confounding}
\begin{columns}
\begin{column}{0.48\textwidth}
\begin{block}{Overfitting}
The model captures noise in the training sample and performs poorly out of sample.
\end{block}
Fix: regularization, cross-validation, honest validation.
\end{column}
\begin{column}{0.48\textwidth}
\begin{alertblock}{Confounding}
The model estimates a relationship that does not represent the causal effect.
\end{alertblock}
Fix: research design, credible variation, identification assumptions.
\end{column}
\end{columns}
\vspace{1em}
\econ{
ML methods primarily address overfitting. Econometric design addresses confounding.
}
\end{frame}

% =========================================================
\begin{frame}{Cross-validation and causality}
Cross-validation answers:
\[
\text{How well does the model predict new observations from a similar data-generating process?}
\]
It does not answer:
\[
\text{Would the outcome change if treatment were assigned differently?}
\]
\begin{alertblock}{Implication}
A cross-validated model can be a strong forecasting tool and a weak causal model.
\end{alertblock}
\end{frame}

% =========================================================
\begin{frame}{Double/debiased ML: preview}
Later in the course, causal ML will use ML to estimate nuisance functions:
\[
m(X)=E[Y|X], \quad e(X)=E[D|X]
\]
Then estimate the causal effect using residualized variables:
\[
Y_i-\widehat m(X_i), \quad D_i-\widehat e(X_i)
\]
\core{
The idea is to use ML flexibly while protecting the target causal parameter from first-stage overfitting.
}
\end{frame}

% =========================================================
\begin{frame}{Orthogonalization intuition}
If we remove the part of $Y$ and $D$ explained by controls $X$, we focus on residual variation:
\[
\tilde Y_i = Y_i-\widehat E[Y_i|X_i]
\]
\[
\tilde D_i = D_i-\widehat E[D_i|X_i]
\]
Then regress:
\[
\tilde Y_i = \theta \tilde D_i + \text{error}
\]
\grad{
This approach is useful when $X$ is high-dimensional and nuisance functions are complex, but identification still requires assumptions such as conditional exogeneity.
}
\end{frame}

% =========================================================
\begin{frame}{Causal forests: preview}
Causal forests aim to estimate:
\[
\tau(x)=E[Y(1)-Y(0)|X=x]
\]
\begin{itemize}
  \item Instead of one average treatment effect, estimate effects that vary by covariates.
  \item Useful for policy targeting and heterogeneity analysis.
  \item Requires careful sample splitting and causal assumptions.
\end{itemize}
\econ{
In the trade-shock paper, one might ask whether import competition harmed some local labor markets more than others.
}
\end{frame}

% =========================================================
\begin{frame}{A simple way to classify your project}
Ask whether the main contribution is:
\begin{enumerate}
  \item \textbf{Prediction:} better forecast or classification.
  \item \textbf{Causal estimation:} more credible or precise effect estimate.
  \item \textbf{Heterogeneity:} who is affected most.
  \item \textbf{Policy learning:} how to allocate treatment.
  \item \textbf{Measurement:} create a new variable from text, images, or high-dimensional data.
\end{enumerate}
\begin{block}{Project advice}
Be explicit. A project can have more than one goal, but each goal needs its own evaluation standard.
\end{block}
\end{frame}

% =========================================================
\begin{frame}{Module 1 coding mindset}
Even in the first module, code should reflect the empirical question.
\begin{itemize}
  \item Create a data dictionary.
  \item Separate outcome, treatment, controls, instruments, fixed effects, and sample restrictions.
  \item Build prediction and causal specifications separately.
  \item Use train/test splits for prediction.
  \item Use design-based diagnostics for causal claims.
  \item Document every transformation.
\end{itemize}
\end{frame}

% =========================================================
\begin{frame}[fragile]{Pseudo-code: prediction workflow}
\begin{verbatim}
1. Define outcome Y and features X.
2. Split data into train and test sets.
3. Fit candidate models on training data.
4. Tune hyperparameters using cross-validation.
5. Evaluate on held-out test data.
6. Interpret predictive performance.
7. Ask whether the model is useful for the policy task.
\end{verbatim}
\begin{alertblock}{Do not conclude}
``Feature $D$ is important, therefore $D$ causes $Y$.''
\end{alertblock}
\end{frame}

% =========================================================
\begin{frame}[fragile]{Pseudo-code: causal workflow}
\begin{verbatim}
1. Define outcome Y, treatment D, controls X, and timing.
2. State the estimand: ATE, ATT, LATE, CATE, etc.
3. Explain the source of treatment variation.
4. Choose the design: IV, RD, DiD, matching, etc.
5. Estimate the effect.
6. Run design diagnostics and robustness checks.
7. Interpret magnitude, mechanism, and limitations.
\end{verbatim}
\begin{block}{Question}
Where, if anywhere, should ML enter this workflow?
\end{block}
\end{frame}

% =========================================================
\begin{frame}{Class exercise: turn each paper into two questions}
For each replication paper, write:
\begin{enumerate}
  \item One prediction question.
  \item One causal question.
  \item One possible ML extension.
  \item One threat to causal interpretation.
\end{enumerate}
\begin{exampleblock}{Example: Hansen}
Prediction: Who is likely to reoffend? \\
Causality: Does crossing the punishment threshold reduce reoffense? \\
ML extension: heterogeneous effects near cutoff. \\
Threat: manipulation or misspecification around cutoff.
\end{exampleblock}
\end{frame}

% =========================================================
\begin{frame}{Common student mistakes}
\begin{enumerate}
  \item Treating high $R^2$ as evidence of causality.
  \item Using post-treatment variables as controls.
  \item Confusing variable importance with causal importance.
  \item Ignoring timing.
  \item Choosing an ML method before defining the estimand.
  \item Reporting model accuracy without economic interpretation.
  \item Reporting causal estimates without discussing identification.
\end{enumerate}
\end{frame}

% =========================================================
\begin{frame}{How to write a good empirical claim}
\begin{block}{Weak claim}
The random forest shows that institutions are important for GDP.
\end{block}
\begin{alertblock}{Better predictive claim}
Institutions are among the strongest predictors of GDP in out-of-sample prediction, conditional on the available feature set.
\end{alertblock}
\begin{exampleblock}{Better causal claim}
Using settler mortality as an instrument, the estimated effect of institutions on GDP is positive under the maintained exclusion restriction.
\end{exampleblock}
\end{frame}

% =========================================================
\begin{frame}{Interpretation: magnitude matters}
A statistically significant coefficient is not enough.
\begin{itemize}
  \item What is the size of the effect?
  \item Is it economically large?
  \item What is the unit of measurement?
  \item How does it compare to baseline levels?
  \item Is the effect local or general?
  \item What are the policy implications?
\end{itemize}
\econ{
In a trade-shock study, a coefficient should be translated into jobs lost, wage changes, or changes in local labor-market outcomes.
}
\end{frame}

% =========================================================
\begin{frame}{External validity}
Even a credible estimate may apply only to a specific setting.
\begin{itemize}
  \item RD estimates are local to the cutoff.
  \item IV estimates may apply to compliers.
  \item Historical persistence estimates may not generalize to modern policy changes.
  \item A prediction model trained in one period may fail in another.
\end{itemize}
\begin{alertblock}{Course habit}
Always ask: credible for whom, where, and when?
\end{alertblock}
\end{frame}

% =========================================================
\begin{frame}{Ethics and policy relevance}
ML and causal inference both raise ethical issues:
\begin{itemize}
  \item targeting can exclude vulnerable groups;
  \item historical variables may encode injustice;
  \item predictions may reinforce existing inequalities;
  \item policy evaluation can affect resource allocation;
  \item transparency matters for public trust.
\end{itemize}
\econ{
Economic interpretation should include distributional consequences, not only average performance.
}
\end{frame}

% =========================================================
\begin{frame}{What a strong Module 1 presentation would do}
A strong paper-method-coding presentation should:
\begin{enumerate}
  \item explain the paper's economic question;
  \item define outcome, treatment, and sample;
  \item explain the identification strategy;
  \item identify the main threat to validity;
  \item show a small reproducible code example;
  \item interpret results economically;
  \item propose one credible ML extension.
\end{enumerate}
\end{frame}

% =========================================================
\begin{frame}{Questions for discussion}
\begin{enumerate}
  \item Can a model be useful for policy even if it is not causal?
  \item When is prediction enough?
  \item When does a policymaker need a causal estimate?
  \item Which replication paper has the strongest identification strategy?
  \item Which replication paper seems most promising for an ML extension?
  \item What are the dangers of using ML in historical or policy settings?
\end{enumerate}
\end{frame}

% =========================================================
\begin{frame}{Module 1 summary}
\begin{itemize}
  \item Prediction and causality are different empirical goals.
  \item Economic questions determine estimands and methods.
  \item Endogeneity and selection bias make causal inference difficult.
  \item Identification is about credible variation, not model complexity.
  \item ML can improve prediction, control selection, heterogeneity analysis, and policy targeting.
  \item ML does not replace economic reasoning or causal design.
  \item The replication papers provide concrete settings where these distinctions matter.
\end{itemize}
\end{frame}

% =========================================================
\begin{frame}{Key formulas to remember}
\begin{block}{Prediction}
\[
\widehat f = \arg\min_f E[L(Y,f(X))]
\]
\end{block}
\begin{block}{Causal effect}
\[
\tau = E[Y(1)-Y(0)]
\]
\end{block}
\begin{block}{Selection bias}
\[
E[Y|D=1]-E[Y|D=0] = ATT + Selection
\]
\end{block}
\begin{block}{Propensity score}
\[
e(X)=P(D=1|X)
\]
\end{block}
\end{frame}

% =========================================================
\begin{frame}{Reading list for Module 1}
\begin{itemize}
  \item Athey, S. (2019). \href{https://www.nber.org/chapters/c14009}{The impact of machine learning on economics}.
  \item Mullainathan, S., and Spiess, J. (2017). Machine learning: an applied econometric approach.
  \item Acemoglu, Johnson, and Robinson (2001). Colonial origins of comparative development.
  \item Nunn (2008). Slave trades and current economic performance.
  \item Hansen (2015). Punishment and deterrence.
  \item Autor, Dorn, and Hanson (2013). The China Syndrome.
\end{itemize}
\smallcite{Additional course papers are used as replication and ML-extension examples.}
\end{frame}

% =========================================================
\begin{frame}{Appendix: Paper-to-project template}
\begin{center}
\scriptsize
\begin{tabular}{p{0.18\textwidth}p{0.15\textwidth}p{0.15\textwidth}p{0.17\textwidth}p{0.22\textwidth}}
\toprule
Paper & Outcome & Treatment / exposure & Design & Possible ML extension \\
\midrule
AJR & Income & Institutions & IV & control robustness, heterogeneity \\
Nunn & Income & Slave trade exposure & historical variation / IV logic & rich controls \\
Goldin et al. & Plan behavior & Subsidy incentives & policy design & classification, anomaly detection \\
Hansen & Recidivism & Punishment threshold & RD & heterogeneous deterrence \\
Anderson & Applications / donations & Athletic success & propensity score & ML propensity, DR estimation \\
ADH & Labor outcomes & Import exposure & shift-share/IV & vulnerability prediction \\
\bottomrule
\end{tabular}
\end{center}
\end{frame}

% =========================================================
\begin{frame}{Appendix: Example project question templates}
\begin{itemize}
  \item Can ML improve prediction of the main outcome relative to the original linear model?
  \item Does flexible covariate adjustment change the estimated treatment effect?
  \item Are treatment effects heterogeneous across economically meaningful groups?
  \item Can a policy-targeting rule based on predicted effects improve welfare?
  \item Does an interpretable ML model reveal mechanisms consistent with the original paper?
\end{itemize}
\begin{alertblock}{Minimum standard}
Every project must separate predictive performance from causal interpretation.
\end{alertblock}
\end{frame}

% =========================================================
\begin{frame}{Appendix: Diagnostic checklist}
\begin{center}
\begin{tabular}{p{0.30\textwidth}p{0.55\textwidth}}
\toprule
If using... & Check... \\
\midrule
Prediction & holdout performance, leakage, calibration, external validity \\
IV & first stage, exclusion restriction, monotonicity, weak instruments \\
RD & manipulation, bandwidth, continuity, local balance \\
DiD & pre-trends, composition changes, spillovers \\
Matching/PS & overlap, balance, sensitivity to unobserved confounding \\
ML causal methods & sample splitting, orthogonality, nuisance performance, assumptions \\
\bottomrule
\end{tabular}
\end{center}
\end{frame}

% =========================================================
\begin{frame}{Appendix: A final mental model}
\begin{center}
\Large
Question $\rightarrow$ Estimand $\rightarrow$ Identification $\rightarrow$ Estimation $\rightarrow$ Validation $\rightarrow$ Interpretation
\end{center}
\vspace{1em}
\core{
Machine learning mostly improves estimation and prediction. Economics determines the question, estimand, identification, and interpretation.
}
\end{frame}

% =========================================================
\begin{frame}{End of Module 1}
\begin{center}
\Huge Questions?
\end{center}
\vspace{1em}
\begin{center}
\large Next: supervised and unsupervised learning in economics.
\end{center}
\end{frame}

\end{document}
